\section{Analiza Wyników i Działania Systemu}
\label{sec:Analiza_wynikow}

\subsection{Metodologia porównania}

Celem eksperymentów jest ilościowe i jakościowe porównanie dwóch podejść do zarządzania cyklem życia modelu uczenia maszynowego \cite{kreuzberger2023mlops}:

\begin{enumerate}
  \item \textbf{Podejście ręczne} --- wszystkie kroki (przygotowanie danych, trening, walidacja, wdrożenie) wykonywane przez człowieka bez automatyzacji.
  \item \textbf{Zautomatyzowane CI/CD (niniejszy projekt)} --- pełna automatyzacja od push danych do wdrożenia modelu w oparciu o GitHub Actions, DVC i MLflow.
\end{enumerate}

Porównanie obejmuje następujące wymiary:
\begin{itemize}
  \item czas nakładu pracy człowieka w każdym podejściu,
  \item jakość modelu (metryki Dice i IoU),
  \item pracochłonność (liczba ręcznych kroków),
  \item ryzyko błędu ludzkiego,
  \item skalowalność procesu.
\end{itemize}

\subsection{Czasy wykonania scenariuszy}

Tabela \ref{tab:czasy} zestawia szacowane czasy nakładu pracy człowieka dla każdego podejścia w scenariuszu dodania nowej partii danych treningowych i wdrożenia ulepszonego modelu. Wartości dla podejścia ręcznego mają charakter szacunkowy i zakładają doświadczonego inżyniera ML znającego narzędzia projektu. Wartości dla podejścia zautomatyzowanego zostały zmierzone na podstawie logów GitHub Actions.

\begin{table}[H]
\centering
\caption{Porównanie nakładu pracy człowieka w obu scenariuszach}
\label{tab:czasy}
\begin{tabular}{|l|r|r|}
\hline
\textbf{Etap} & \textbf{Ręczny} & \textbf{Zautomatyzowane CI/CD} \\
\hline
Przygotowanie i upload danych & $\sim$30 min & $\sim$10 min \\
Walidacja danych & $\sim$30 min & $\sim$10 min (poprawki złych danych) \\
Trening modelu & $\sim$3,5 h & automatyczny \\
Ocena jakości & $\sim$15 min & automatyczna \\
Wdrożenie & $\sim$45 min & automatyczne \\
\hline
\textbf{Nakład człowieka łącznie} & $\sim$5 h & $\sim$20 min \\
\textbf{Całkowity czas pipeline'u} & $\sim$5 h & $\sim$4 h \\
\hline
\end{tabular}
\end{table}

Kluczową obserwacją jest rozróżnienie między \textit{nakładem pracy człowieka} a \textit{całkowitym czasem realizacji}. W podejściu ręcznym są one tożsame --- człowiek jest zaangażowany przez cały czas trwania procesu. W podejściu zautomatyzowanym inżynier wykonuje jedynie komendę \texttt{git push} (poniżej jednej minuty), po czym pipeline działa bez jego udziału. Całkowity czas działania pipeline'u, od push do wdrożenia, można odczytać z zakładki \textit{Actions} w repozytorium GitHub.

\subsection{Metryki modelu}

Tabela \ref{tab:metryki} przedstawia metryki modelu dla kolejnych wersji wytrenowanych w ramach projektu. Baseline pochodzi z repozytorium referencyjnego \textit{Water-Meters-Segmentation} \cite{kaggle_water_meters}. Metryki są rejestrowane automatycznie przez MLflow \cite{zaharia2018mlflow} podczas każdego treningu.

\begin{table}[H]
\centering
\caption{Metryki jakości modelu dla kolejnych wersji}
\label{tab:metryki}
\begin{tabular}{|l|c|c|r|}
\hline
\textbf{Wersja modelu} & \textbf{Dice} & \textbf{IoU} & \textbf{Hausdorff [px]} \\
\hline
v1  & 0{,}0042 & 0{,}0021 & 358{,}61 \\
v2  & 0{,}0741 & 0{,}0388 & 148{,}56 \\
v3  & 0{,}4710 & 0{,}4060 & 181{,}28 \\
v4  & 0{,}4825 & 0{,}4518 &  94{,}33 \\
v5  & 0{,}5998 & 0{,}5508 & 226{,}52 \\
v6 (Production)  & 0{,}8134 & 0{,}7409 &  60{,}00 \\
\hline
\end{tabular}
\end{table}

Postęp metryk ilustruje działanie bramki jakości: każda kolejna wersja trafiająca do etapu Production musiała uzyskać wyższy Dice i IoU od poprzedniej. Metrykę Hausdorff Distance obliczono na zbiorze testowym jako odległość Hausdorffa między maską przewidywaną a maską referencyjną; mniejsza wartość oznacza lepsze dopasowanie konturów. Widoczna poprawa z 358 px (v1) do 60 px (v6) potwierdza postępującą naukę modelu.

\begin{figure}[H]
  \centering
  \includegraphics[width=0.85\linewidth]{img/07_analiza_wyników_i_dzialanie_systemu/mlflow_modes_metrics.png}
  \caption{Metryki kolejnych wersji modelu w interfejsie MLflow.}
  \label{fig:mlflow-metrics}
\end{figure}

Poniżej przedstawiono wykresy przebiegu treningu najlepszej wersji modelu (v6) na przestrzeni 50 epok. Wykresy generowane są automatycznie przez skrypt \texttt{train.py} i archiwizowane w MLflow jako artefakty każdego eksperymentu.

\begin{figure}[H]
  \centering
  \begin{subfigure}[b]{0.48\linewidth}
    \includegraphics[width=\linewidth]{img/07_analiza_wyników_i_dzialanie_systemu/plot_dice.png}
    \caption{Współczynnik Dice}
    \label{fig:plot-dice}
  \end{subfigure}
  \hfill
  \begin{subfigure}[b]{0.48\linewidth}
    \includegraphics[width=\linewidth]{img/07_analiza_wyników_i_dzialanie_systemu/plot_iou.png}
    \caption{Współczynnik IoU}
    \label{fig:plot-iou}
  \end{subfigure}

  \vspace{0.8em}

  \begin{subfigure}[b]{0.48\linewidth}
    \includegraphics[width=\linewidth]{img/07_analiza_wyników_i_dzialanie_systemu/plot_loss.png}
    \caption{Funkcja straty}
    \label{fig:plot-loss}
  \end{subfigure}
  \hfill
  \begin{subfigure}[b]{0.48\linewidth}
    \includegraphics[width=\linewidth]{img/07_analiza_wyników_i_dzialanie_systemu/plot_accuracy.png}
    \caption{Dokładność pikselowa}
    \label{fig:plot-accuracy}
  \end{subfigure}

  \caption{Przebieg treningu modelu v6 na przestrzeni 50 epok (zbiory: treningowy, walidacyjny, testowy).}
  \label{fig:training-curves}
\end{figure}

Rysunek \ref{fig:mlflow-pred-2} przedstawia przykładowy wynik inferencji modelu v6: obraz wejściowy, maskę referencyjną (ground truth) oraz maskę przewidywaną. Widoczna zbieżność konturów wodomierza potwierdza praktyczną użyteczność modelu.

\begin{figure}[H]
  \centering
  \includegraphics[width=0.85\linewidth]{img/07_analiza_wyników_i_dzialanie_systemu/plot_pred_2.png}
  \caption{Przykład predykcji modelu v6: obraz wejściowy (lewo), maska referencyjna (środek), maska przewidywana (prawo).}
  \label{fig:mlflow-pred-2}
\end{figure}

\subsection{Koszty infrastruktury AWS}

\begin{table}[H]
\centering
\caption{Szacowane koszty infrastruktury AWS}
\label{tab:koszty}
\begin{tabular}{|l|r|l|}
\hline
\textbf{Zasób} & \textbf{Koszt} & \textbf{Uwagi} \\
\hline
EC2 t3.large (trening i serwowanie) & \$0{,}0832/h & Uruchamiana tylko podczas treningu (stop/start) \\
S3 (artefakty DVC i MLflow) & $\sim$\$0{,}02/mies. & Kilka GB artefaktów \\
ECR (obrazy Docker) & $\sim$\$0{,}05/mies. & Jeden obraz $\approx$500\,MB \\
\hline
\textbf{Łącznie (projekt)} & $\sim$\$20 & Za cały okres realizacji \\
\hline
\end{tabular}
\end{table}

Efektywność kosztową zapewnia architektura stop/start instancji EC2: instancja uruchamiana jest wyłącznie na czas treningu (job \texttt{start-infra}) i zatrzymywana bezpośrednio po jego zakończeniu (job \texttt{stop-infra}, uruchamiany zawsze niezależnie od wyniku). Przy stawce \$0{,}0832/h i czasie treningu 3,5 h, koszt pojedynczej sesji treningowej wynosi zaledwie $\sim$\$0{,}29.

\subsection{Skalowalność wysiłku i kosztów}

Przedstawione do tej pory wyniki dotyczą zbioru 359 obrazów. W rzeczywistych zastosowaniach przemysłowych zbiory danych rosną wraz z eksploatacją systemu. Poniżej przeanalizowano, jak oba podejścia zachowują się przy zwiększaniu rozmiaru danych.

Czas treningu skaluje się w przybliżeniu liniowo z liczbą obrazów, ponieważ każda epoka przetwarza cały zbiór treningowy. Przy obserwowanym czasie $\sim$3,5 h dla 359 obrazów, szacunki dla większych zbiorów przedstawia tabela~\ref{tab:skalowanie}.

\begin{table}[H]
\centering
\caption{Szacowane czasy treningu i koszty w zależności od rozmiaru zbioru}
\label{tab:skalowanie}
\begin{tabular}{|r|r|r|r|r|}
\hline
\textbf{Obrazy} & \textbf{Czas treningu} & \textbf{Koszt EC2} & \textbf{Nakład ręczny*} & \textbf{Koszt pracownika**} \\
\hline
359   & $\sim$3{,}5 h &  $\sim$\$0{,}29 & $\sim$5 h   & $\sim$150 PLN \\
700   & $\sim$7 h    &  $\sim$\$0{,}58 & $\sim$7 h   & $\sim$210 PLN \\
1500  & $\sim$15 h   &  $\sim$\$1{,}25 & $\sim$10 h  & $\sim$300 PLN \\
5000  & $\sim$50 h   &  $\sim$\$4{,}16 & $\sim$20 h  & $\sim$600 PLN \\
\hline
\end{tabular}
\end{table}

\noindent\small{*Nakład ręczny uwzględnia czas oczekiwania na wynik treningu, ocenę modelu i ręczne wdrożenie.}\\
\noindent\small{**Szacunkowo przy stawce 30~PLN/h (minimalne wynagrodzenie w 2025 r.); przy wyższych stawkach dysproporcja jest jeszcze większa.}

\vspace{0.5em}

Z tabeli wynikają dwie kluczowe obserwacje:

\begin{enumerate}
  \item \textbf{Koszt AWS rośnie wolno i pozostaje pomijalny.} Przy 5000 obrazach koszt jednego treningu wynosi $\sim$\$4 (ok.~17~PLN). Nawet przy cyklach miesięcznych roczny koszt obliczeniowy to zaledwie kilkadziesiąt złotych.

  \item \textbf{Koszt pracownika rośnie szybko i zdominuje budżet.} Przy ręcznym podejściu wzrost zbioru danych przekłada się bezpośrednio na więcej czasu oczekiwania, więcej kroków manualnych i wyższe ryzyko pomyłki. Przy 5000 obrazach jeden cykl wdrożenia angażuje pracownika przez $\sim$20 h, co przy miesięcznych aktualizacjach daje roczny koszt rzędu 7200~PLN na samą obsługę modelu.
\end{enumerate}

Co istotne, w podejściu zautomatyzowanym wzrost rozmiaru zbioru nie zwiększa nakładu pracy człowieka --- inżynier wykonuje tę samą jedną komendę \texttt{git push} niezależnie od tego, czy trenuje model na 359 czy na 5000 obrazach. Cały ciężar obliczeniowy przejmuje infrastruktura AWS.

Warto jednak zauważyć, że podejście ręczne może zachować uzasadnienie ekonomiczne w wąskim scenariuszu: bardzo rzadkie aktualizacje modelu (raz na rok lub rzadziej) przy małym zbiorze danych i braku wymagań co do powtarzalności. W takim przypadku jednorazowy koszt konfiguracji automatyzacji może nie zwrócić się przed końcem projektu. Niemniej wraz ze wzrostem częstotliwości aktualizacji lub rozmiaru danych próg opłacalności automatyzacji jest przekraczany coraz szybciej.

\subsection{Ocena jakościowa}

\subsubsection{Łatwość wdrożenia}

Konfiguracja systemu wymagała jednorazowego nakładu pracy obejmującego: przygotowanie infrastruktury Terraform (VPC, EC2, S3, ECR), konfigurację k3s i Helm, integrację MLflow z S3, napisanie skryptów walidacji danych i bramki jakości oraz przygotowanie workflow GitHub Actions. Jednorazowy koszt konfiguracji wyniósł kilka tygodni pracy. Po jego poniesieniu każde kolejne wdrożenie nowej wersji modelu wymaga jedynie wykonania \texttt{git push}.

Podejście ręczne nie wymaga nakładu konfiguracyjnego, ale generuje cykliczny nakład pracy przy każdym wdrożeniu. Przy założeniu jednego wdrożenia miesięcznie i nakładzie $\sim$5 h na cykl automatyzacja staje się opłacalna po kilkunastu cyklach, nie uwzględniając ryzyka błędu ludzkiego, które rośnie z liczbą powtórzeń.

\subsubsection{Stabilność systemu}

W trakcie testów zaobserwowano następujące zachowania systemu w sytuacjach wyjątkowych:
\begin{itemize}
  \item \textbf{Niepoprawne dane} --- walidator odrzucał dane z komunikatem i nie uruchamiał treningu,
  \item \textbf{Model gorszy od baseline} --- bramka jakości odrzucała model bez tworzenia PR,
  \item \textbf{Awaria instancji EC2} --- job \texttt{stop-infra} uruchamiał się zawsze (\texttt{if: always()}), zapobiegając wyciekom kosztów.
\end{itemize}

\subsubsection{Skalowalność}

Zastosowane podejście ma następujące ograniczenia skalowalności:
\begin{itemize}
  \item single-node k3s --- brak wysokiej dostępności,
  \item jeden runner EC2 --- brak równoległych treningów,
  \item model proof of concept --- nie testowano przy zbiorach powyżej 1000 obrazów.
\end{itemize}
